% Book pp. 58-62 (PDF pp. 59-63)
\section{Recursive Sequences}

A recursive sequence is a sequence where each term is defined using previous
terms.

For example, for an Arithmetic Sequence:

First term $a_1 = 2$

Recursive formula: $a_n = a_{n-1} + 3$
\begin{enumerate}
  \item Start with $a_1 = 2$
  \item To find $a_2$, add 3 to $a_1$ : $a_2 = 2 + 3 = 5$
  \item To find $a_3$, add 3 to $a_2$ : $a_3 = 5 + 3 = 8$
  \item To find $a_4$, add 3 to $a_3$ : $a_4 = 8 + 3 = 11$

  Sequence: 2, 5, 8, 11, 14, 17\ldots
\end{enumerate}

For example, for a Geometric Sequence:

First Term: $a_1 = 2$

Recursive formula: $a_n = 2 \times a_{n-1}$
\begin{enumerate}
  \item Start with $a_1 = 2$
  \item To find $a_2$, multiply $a_1$ \textit{by} $2 \to a_2 = 2 \times 2 = 4$
  \item To find $a_3$, multiply $a_2$ \textit{by} $2 \to a_3 = 2 \times 4 = 8$
  \item To find $a_4$, multiply $a_3$ \textit{by} $2 \to a_4 = 2 \times 8 = 16$

  Sequence: 2, 4, 8, 16, 32, 64, \ldots
\end{enumerate}

In both examples;
\begin{itemize}
  \item The first term $(a_1)$ is given
  \item The recursive formula defines how to find the next term $(a_n)$ using the
  previous term $(a_{n-1})$
\end{itemize}

\subsection{Recurrent decimals}

Recurrent decimals are repeated decimals. It can be seen that every recurrent
decimal represents a rational number. Every rational number can be represented
by a terminal decimal or a recurrent decimal. If the denominator of a rational
number (written in its simplest form) has prime factors, other than 2 and 5, then
the rational number is recognised as a recurrent decimal.

For example; $\frac{5}{9} = 0.5555\ldots$ \textit{or} $0.\dot{5}$,
$\frac{2}{11} = 0.181818\ldots$ \textit{or} $0.\dot{1}\dot{8}$,
$\frac{8}{15} = 0.53333\ldots$ \textit{or} $0.5\dot{3}$

\begin{example}{Example 2.7}
Write the recurring decimals as a series. Identify the common ratio, the first term
and an explicit formula for the $n$th term:
\begin{enumerate}
  \item[a.] 0.222222\ldots
  \item[b.] 0.323232\ldots
  \item[c.] 5.414141\ldots
\end{enumerate}

\solution
\begin{enumerate}
  \item[a.] $0.2222... = 0.2 + 0.02 + 0.002 + 0.0002 + \cdots$
  \begin{align*}
    &= 0.2 + 0.2(0.1) + 0.2(0.01) + 0.2(0.001) + \cdots \\
    &= 0.2 + 0.2(0.1) + 0.2(0.1)^2 + 0.2(0.1)^3 + \cdots
  \end{align*}
  The first term is 0.2, and the common ratio between each successive term is
  0.1. The formula for the $n$th term will be $U_n = 0.2\left(\frac{1}{10}\right)^{n-1}$

  \item[b.] $0.323232\ldots = 0.32 + 0.0032 + 0.000032 + 0.00000032 + \cdots$
  \begin{align*}
    &= 0.32 + 0.32(0.01) + 0.32(0.0001) + 0.32(0.000001) + \cdots \\
    &= 0.32 + 0.32(0.01) + 0.32(0.01)^2 + 0.32(0.01)^3 + \cdots
  \end{align*}
  The first term is 0.32, and the common ratio between each successive term
  is 0.01. The formula for the $n$th term will be $U_n = 0.32\left(\frac{1}{100}\right)^{n-1}$

  \item[c.] $5.414141\ldots = 5 + 0.41 + 0.0041 + 0.000041 + 0.00000041 + \cdots$
  \[
    = 5 + 0.41 + 0.41(0.01) + 0.41(0.1)^2 + 0.41(0.1)^3 + \cdots
  \]
  From the second term (0.41), a geometric sequence with first term 0.41, and
  the common ratio 0.01 can be observed. The formula for the $n$th term will
  be

  For 0.414141\ldots, $U_n = 0.41\left(\frac{1}{100}\right)^{n-1}$

  Thus, $5.414141\ldots = 5 + 0.41\left(\frac{1}{100}\right)^{n-1}$ for $n = 1, 2, 3, \ldots$
  \booknote{the book prints $0.41(0.1)^2 + 0.41(0.1)^3$; with common ratio $0.01$ these
  should read $0.41(0.01)^2 + 0.41(0.01)^3$.}
\end{enumerate}
\end{example}

\begin{example}{Example 2.8}
Write 0.9999\ldots as an exponential series and show that in the limit, the sum of
series equals one.

\solution
$0.9999\ldots = 0.9 + 0.09 + 0.009 + 0.0009 + \ldots$

This is a G.P. with first term, $a = 0.9$ and common ratio, $r = \frac{0.09}{0.9} = 0.1$

The series is an infinite series with sum; $S_\infty = \frac{a}{1 - r} = \frac{0.9}{1 - 0.1}
= \frac{0.9}{0.1} = 1$
\booknote{$1 - 0.1 = 0.9$, so the last step should read $\frac{0.9}{0.9} = 1$; the book
prints $\frac{0.9}{0.1}$.}
\end{example}

\begin{example}{Example 2.9}
Express $0.1\dot{6}$ as an infinite geometric series and find the sum of the geometric
series.

\solution
\begin{align*}
  0.1\dot{6} &= 0.166666\ldots \\
  &= 0.1 + 0.06 + 0.006 + 0.0006 + \ldots \\
  &= \frac{1}{10} + \frac{6}{100} + \frac{6}{1000} + \frac{6}{10000} + \ldots
\end{align*}
From the series, the term after $\frac{1}{10}$ from geometric series with $= \frac{6}{100}$,
$r = \frac{6}{100} \div \frac{6}{1000} = \frac{1}{10}$

Sum of the geometric series:
\begin{align*}
  0.0\dot{6} = S_\infty &= \frac{\frac{6}{100}}{1 - \frac{1}{10}} \\
  &= \frac{\frac{6}{100}}{\frac{9}{10}} \\
  &= \frac{1}{15}
\end{align*}
$\therefore 0.1\dot{6} = \frac{1}{10} + \frac{1}{15} = \frac{1}{6}$
\booknote{the common ratio is $\frac{6}{1000} \div \frac{6}{100} = \frac{1}{10}$; the book
prints the division the other way round.}
\end{example}

\begin{example}{Example 2.10}
If a sequence $U_1, U_2, U_3\ldots$ is define by the relation $U_{n+1} = 3 + U_n$ for
$n \ge 1$ and $U_1 = 1$, find;
\begin{enumerate}
  \item[a.] $U_2, U_3$ and $U_4$
  \item[b.] a formula for $U_n$ in terms of $n$.
  \item[c.] the sum of the 1\textit{st} $n$ terms
\end{enumerate}

\solution
\begin{enumerate}
  \item[a.] Using the relation
  \begin{align*}
    U_{n+1} &= 3 + U_n \\
    \text{If } n &= 1 \\
    U_2 &= 3 + U_1 = 3 + 1 = 4 \\
    U_3 &= 3 + U_2 = 3 + 4 = 7 \\
    U_4 &= 3 + U_3 = 3 + 7 = 10
  \end{align*}
  The sequence is $1, 4, 7, 10, \ldots$, so it is a linear sequence
  \item[b.] $U_n = a + (n-1)d$ and $a = 1, d = 3$
  \begin{align*}
    U_n &= 1 + (n-1)3 \\
    U_n &= 1 + 3n - 3 \\
    U_n &= 3n - 2
  \end{align*}
  \item[c.] The sum of $n$ terms of $1, 4, 7, 10, \ldots$
  \begin{align*}
    S_n &= \frac{n}{2}[2a + (n-1)d] \\
    S_n &= \frac{n}{2}[2(1) + (n-1)3] \\
    S_n &= \frac{n}{2}[2 + 3n - 3] \\
    S_n &= \frac{n}{2}[3n - 1]
  \end{align*}
\end{enumerate}
\end{example}
